% !TEX root = ../main.tex
\chapter{项目组织与管理计划}

\section{项目组织结构}

\subsection{组织架构概述}

本项目采用\textbf{矩阵式组织结构}，结合敏捷开发理念，建立高效的小型团队。项目组织架构如图~\ref{fig:org_structure}所示。

\begin{figure}[htbp]
    \centering
    \begin{tikzpicture}[
            node distance=1.2cm and 2cm,
            box/.style={rectangle, draw, rounded corners, minimum width=3cm, minimum height=1cm, align=center, font=\small},
            leader/.style={box, fill=blue!30},
            member/.style={box, fill=green!20},
            advisor/.style={box, fill=orange!20},
            arrow/.style={->, >=stealth, thick}
        ]
        % 指导老师
        \node[advisor] (advisor) {指导老师\\殷凯};

        % 项目负责人
        \node[leader, below=of advisor] (pm) {项目负责人\\许子祺};

        % 团队成员
        \node[member, below left=1.5cm and 1cm of pm] (dev1) {后端开发\\许子祺};
        \node[member, below right=1.5cm and 1cm of pm] (dev2) {前端开发\\莫文涛};

        % 职责标注
        \node[below=0.3cm of dev1, font=\footnotesize, text=gray, align=center] {Flask API\\ML服务\\Gurobi优化};
        \node[below=0.3cm of dev2, font=\footnotesize, text=gray, align=center] {Flutter Web\\数据可视化\\用户交互};

        % 箭头
        \draw[arrow] (advisor) -- node[right, font=\footnotesize] {指导} (pm);
        \draw[arrow] (pm) -- (dev1);
        \draw[arrow] (pm) -- (dev2);

        % 协作关系
        \draw[<->, dashed, thick, gray] (dev1) -- node[below, font=\footnotesize] {API协作} (dev2);
    \end{tikzpicture}
    \caption{项目组织结构图}
    \label{fig:org_structure}
\end{figure}

\subsection{角色定义}

表~\ref{tab:roles}定义了各角色的职责范围。

\begin{table}[htbp]
    \centering
    \small
    \caption{项目角色定义}
    \label{tab:roles}
    \begin{tabular}{p{2cm}p{2cm}p{9cm}}
        \toprule
        \textbf{角色} & \textbf{人员} & \textbf{职责描述}              \\
        \midrule
        指导老师        & 殷凯          & 技术指导、文档审核、协调Gurobi学术许可     \\
        项目负责人       & 许子祺         & 项目规划、架构设计、后端开发、文档统筹        \\
        前端开发        & 莫文涛         & Flutter Web开发、数据可视化、UI交互设计 \\
        \bottomrule
    \end{tabular}
\end{table}

\section{人员分工职责}

\subsection{工作分解结构 (WBS)}

项目的工作分解结构如图~\ref{fig:wbs}所示，将整体工作划分为5个主要工作包。

\begin{figure}[htbp]
    \centering
    \begin{tikzpicture}[
            level 1/.style={sibling distance=2.8cm, level distance=1.6cm},
            level 2/.style={sibling distance=1.4cm, level distance=1.3cm},
            every node/.style={rectangle, draw, rounded corners, minimum width=1.2cm, minimum height=0.6cm, align=center, font=\tiny}
        ]
        \node[fill=blue!20, font=\footnotesize] {智能能源\\管理平台}
        child {node[fill=green!20] {需求分析\\WP1}
                child {node {调研}}
                child {node {文档}}
            }
        child {node[fill=green!20] {系统设计\\WP2}
                child {node {架构}}
                child {node {接口}}
            }
        child {node[fill=orange!20] {开发实现\\WP3}
                child {node {后端}}
                child {node {前端}}
            }
        child {node[fill=yellow!30] {测试部署\\WP4}
                child {node {测试}}
                child {node {部署}}
            }
        child {node[fill=purple!20] {文档编写\\WP5}
                child {node {技术}}
                child {node {手册}}
            };
    \end{tikzpicture}
    \caption{工作分解结构 (WBS)}
    \label{fig:wbs}
\end{figure}


\subsection{详细任务分配}

表~\ref{tab:task_assignment}展示了各工作包下的具体任务分配情况。

\begin{table}[H]
    \centering
    \caption{详细任务分配表}
    \label{tab:task_assignment}
    \begin{tabular}{clcc}
        \toprule
        \textbf{工作包}                    & \textbf{任务}   & \textbf{负责人} & \textbf{工时(人天)} \\
        \midrule
        \multirow{3}{*}{WP1: 需求分析}

                                        & 功能需求文档        & 许子祺          & 2               \\
                                        & 非功能需求定义       & 许子祺          & 1               \\
        \midrule
        \multirow{4}{*}{WP2: 系统设计}
                                        & 系统架构设计        & 许子祺          & 3               \\
                                        & 数据库/存储设计      & 许子祺          & 2               \\
                                        & API接口设计       & 许子祺          & 2               \\
                                        & UI/UX设计       & 莫文涛          & 3               \\
        \midrule
        \multirow{6}{*}{WP3: 开发实现}
                                        & Flask后端框架搭建   & 许子祺          & 3               \\
                                        & ML预测服务开发      & 许子祺          & 5               \\
                                        & Gurobi优化服务开发  & 许子祺          & 4               \\
                                        & Flutter前端框架搭建 & 莫文涛          & 3               \\
                                        & 数据可视化组件开发     & 莫文涛          & 5               \\
                                        & 用户交互功能实现      & 莫文涛          & 4               \\
        \midrule
        \multirow{3}{*}{WP4: 测试部署}
                                        & 单元测试编写        & 许子祺、         & 3               \\
                                        & 集成测试          & 许子祺          & 2               \\
                                        & GAE云端部署       & 许子祺          & 2               \\
        \midrule
        \multirow{3}{*}{WP5: 文档编写}
                                        & 可行性研究报告       & 许子祺          & 4               \\
                                        & 项目开发计划书       & 许子祺          & 3               \\
                                        & 技术报告          & 莫文涛          & 4               \\
        \midrule
        \multicolumn{2}{r}{\textbf{总计}} &               & \textbf{57}                    \\
        \bottomrule
    \end{tabular}
\end{table}

\section{进度安排与里程碑}

\subsection{项目时间轴}

本项目计划于\textbf{2024年10月}启动，\textbf{2025年1月}完成，总工期约4个月。项目进度甘特图如图~\ref{fig:gantt}所示。

\begin{figure}[htbp]
    \centering
    \begin{tikzpicture}[x=0.8cm, y=0.6cm]
        % 时间轴
        \draw[thick, ->] (0, 0) -- (18, 0) node[right] {时间};
        \foreach \x/\month in {0/10月, 4/11月, 8/12月, 12/1月, 16/2月} {
                \draw (\x, -0.2) -- (\x, 0.2);
                \node[below] at (\x, -0.3) {\footnotesize \month};
            }

        % 任务条
        \def\task#1#2#3#4#5{
            \fill[#5] (#2, #1) rectangle (#3, #1+0.5);
            \node[right, font=\footnotesize] at (-0.1, #1+0.25) {#4};
        }

        \task{6}{0}{3}{需求分析}{blue!40}
        \task{5}{2}{6}{系统设计}{green!40}
        \task{4}{5}{12}{后端开发}{orange!40}
        \task{3}{6}{13}{前端开发}{yellow!60}
        \task{2}{11}{14}{测试验证}{purple!40}
        \task{1}{13}{16}{部署上线}{red!40}

        % 里程碑
        \foreach \x/\label in {3/M1, 6/M2, 12/M3, 14/M4, 16/M5} {
                \fill[red] (\x, 7) -- (\x-0.3, 7.5) -- (\x+0.3, 7.5) -- cycle;
                \node[above, font=\footnotesize\bfseries] at (\x, 7.5) {\label};
            }

        % 图例
        \node[font=\tiny, right, align=left] at (0, -1.5) {M1:需求冻结 M2:设计评审 M3:开发完成 M4:测试通过 M5:发布};
    \end{tikzpicture}
    \caption{项目进度甘特图}
    \label{fig:gantt}
\end{figure}

\subsection{里程碑定义}

表~\ref{tab:milestones}定义了项目的关键里程碑及其验收标准。

\begin{table}[htbp]
    \centering
    \caption{项目里程碑定义}
    \label{tab:milestones}
    \begin{tabular}{clp{5cm}l}
        \toprule
        \textbf{里程碑} & \textbf{名称} & \textbf{验收标准}            & \textbf{计划日期} \\
        \midrule
        M1           & 需求冻结        & 需求文档通过评审，无重大变更           & 2024-10-31    \\
        M2           & 设计评审        & 架构设计文档通过评审，API接口定义完成     & 2024-11-15    \\
        M3           & 开发完成        & 核心功能开发完成，通过冒烟测试          & 2024-12-31    \\
        M4           & 测试通过        & 所有测试用例通过，缺陷关闭率$\geq$95\% & 2025-01-15    \\
        M5           & 正式发布        & 系统部署上线，用户验收通过            & 2025-01-25    \\
        \bottomrule
    \end{tabular}
\end{table}

\section{风险评估与管理}

\subsection{风险识别}

根据项目特点，识别出以下主要风险，如表~\ref{tab:risks}所示。

\begin{table}[htbp]
    \centering
    \caption{项目风险评估矩阵}
    \label{tab:risks}
    \begin{tabular}{p{0.8cm}p{3.5cm}cccp{4cm}}
        \toprule
        \textbf{ID} & \textbf{风险描述} & \textbf{概率} & \textbf{影响} & \textbf{等级}            & \textbf{应对措施} \\
        \midrule
        R1          & Gurobi许可证配置失败 & 中           & 高           & \cellcolor{orange!30}高 &
        提前申请学术许可；准备开源替代方案(CBC/GLPK)                                                                      \\
        R2          & GAE部署超时/配额限制  & 中           & 中           & \cellcolor{yellow!30}中 &
        优化代码启动时间；使用懒加载；申请配额提升                                                                            \\
        R3          & ML模型预测精度不达标   & 低           & 高           & \cellcolor{yellow!30}中 &
        增加特征工程；尝试多模型集成；扩充训练数据                                                                            \\
        R4          & 团队成员时间冲突      & 中           & 中           & \cellcolor{yellow!30}中 &
        建立缓冲时间；灵活任务调配；每日站会同步                                                                             \\
        R5          & API接口变更频繁     & 低           & 低           & \cellcolor{green!30}低  &
        冻结接口设计；版本控制；充分设计评审                                                                               \\
        \bottomrule
    \end{tabular}
\end{table}

\subsection{风险应对策略}

针对高风险项，制定详细的应对计划：

\paragraph{R1: Gurobi许可证配置}
\begin{enumerate}
    \item \textbf{预防措施}：项目启动第一周完成Gurobi学术许可申请；
    \item \textbf{缓解措施}：配置WLS（Web License Service）支持云端部署；
    \item \textbf{应急措施}：若许可证失效，切换至开源求解器CBC，牺牲部分求解效率。
\end{enumerate}

\paragraph{R2: GAE部署问题}
\begin{enumerate}
    \item \textbf{预防措施}：采用懒加载策略，延迟初始化重型依赖（如Gurobi、ML模型）；
    \item \textbf{缓解措施}：使用F4实例类型，增加启动超时时间至10分钟；
    \item \textbf{应急措施}：切换至Cloud Run或GKE部署方案。
\end{enumerate}

\section{沟通管理计划}

\subsection{沟通渠道}

表~\ref{tab:communication}定义了项目的沟通渠道和频率。

\begin{table}[htbp]
    \centering
    \caption{沟通管理矩阵}
    \label{tab:communication}
    \begin{tabular}{llll}
        \toprule
        \textbf{沟通类型} & \textbf{频率} & \textbf{参与人} & \textbf{工具} \\
        \midrule
        每日站会          & 每日10:00     & 全体成员         & 微信群         \\
        周进度汇报         & 每周五         & 全体成员 + 指导老师  & 腾讯会议        \\
        代码评审          & 按需          & 开发人员         & GitHub PR   \\
        里程碑评审         & 里程碑节点       & 全体成员 + 指导老师  & 线下会议        \\
        问题上报          & 实时          & 相关人员         & 微信/Issues   \\
        \bottomrule
    \end{tabular}
\end{table}

\subsection{会议制度}

\begin{itemize}
    \item \textbf{每日站会}：时长10分钟，每人回答三个问题：昨天做了什么？今天计划做什么？有什么阻碍？
    \item \textbf{周进度汇报}：时长30分钟，展示本周成果，讨论下周计划，同步风险和问题。
    \item \textbf{里程碑评审}：时长1小时，演示阶段成果，收集反馈，确认下阶段目标。
\end{itemize}